\documentclass{beamer}


%variables
\newcommand{\hhulogo}{images/hhulogo}

\usetheme{CambridgeUS}
\usecolortheme{dolphin}
\usefonttheme{structurebold}
\useinnertheme{rounded}
\useoutertheme{miniframes,smoothbars}

\usepackage{pgfpages}
%\setbeameroption{show notes on second screen}

\setbeamertemplate{navigation symbols}{}
\setbeamertemplate{itemize subitem}[triangle]

\usepackage{ulem}

\usepackage[ngerman]{babel}
\usepackage[latin1]{inputenc}
\usepackage[T1]{fontenc}
\usepackage{lmodern}

\usepackage{color} % Farben
\usepackage{microtype} % spezielle PDF-mikrotypographische Optimierungen
\usepackage{graphicx} % Bilder
\usepackage{multimedia}

%better tabulars
\usepackage{array} % Better tabulars
\usepackage{booktabs}
\newcommand{\PreserveBackslash}[1]{\let\temp=\\#1\let\\=\temp}
\usepackage[NewCommands]{ragged2e} % Better raggedleft formating
\newcolumntype{v}[1]{>{\PreserveBackslash{\raggedright}\hspace{0pt}}p{#1}}
\renewcommand{\arraystretch}{1.20}
\usepackage{multirow}

\usepackage{minted}
\definecolor{bg}{rgb}{0.95,0.95,0.95}

\title[Versionierung mit GIT]{Beispielhafte Einf�hrung in die Versionierung mit GIT im Vergleich zu SVN}
\author{Tobias Witt} 
\institute[HHU]{Institut f�r Informatik\\Programmierpraktikum Sommersemester 2012}
\date[20.03.2012]{20. M�rz 2012}
\logo{\pgfimage[width=2cm]{\hhulogo}}

\begin{document}
%
\begin{frame}[plain]
  \begin{center}
    \includegraphics[width=4cm]{\hhulogo}
  \end{center}
  \vspace{-1cm}
  \titlepage
\end{frame}
%
\section*{Gliederung}
\setcounter{framenumber}{0}
\begin{frame}
  \frametitle{Gliederung}
  \tableofcontents
\end{frame}
%
\section{Einf�hrung in GIT}
\begin{frame}
  \frametitle{Einf�hrung in GIT}
  \tableofcontents[currentsection]
\end{frame}
\subsection{Dezentrale Versionierung}
\begin{frame}
  \frametitle{Dezentrale Versionierung}
  \begin{itemize}
    \item<1->Kopie des Repositories in jeder Arbeitskopie enthalten
      \begin{itemize}
        \item In: \texttt{<project-path>/.git}
        \item Arbeitskopie und Repository synonym
      \end{itemize}
    \item<2->Import und Export zwischen Repositories
      \begin{itemize}
          \item zentrale Repositories zur Zusammenarbeit m�glich
          \item Unterscheidung zwischen lokalen und entfernten Repositories
      \end{itemize}
    \item <3->Fast s�mtliche Versionierungsarbeit im lokalen Repository
      \begin{itemize}
          \item Wesentlich h�here Arbeitsgeschwindigkeit
      \end{itemize}
    \item <4->Freiheit f�r die Entwickler
  \end{itemize}
\end{frame}
\subsection[Trees, Commits \& Blobs]{Aufbau und Bestandteile von GIT-Repositories}
\begin{frame}
  \frametitle{Trees, Commits \& Blobs}
  \begin{itemize}
    \item Jede Version in Objektdatenbank in \texttt{.git/objects}
    \item <2->Blob-Objekte
    \item <3->Tree-Objekte
    \item <4->Commit-Objekte
  \end{itemize}
\end{frame}
\begin{frame}
  \frametitle{Blob-Objekte}
  \begin{itemize}
    \item Repr�sentieren Inhalte (von Dateien)
    \item Dargestellt durch einen SHA1-Hashwert (u.A. abh�ngig vom Inhalt der Datei)
    \item Gleicher Datei-Inhalt $\Rightarrow$ Wiederverwendung des Blob-Objekts
    \item Kein Blob-Objekt, keine Versionierung
      \begin{itemize}
        \item Leere Verzeichnisse werden nicht versioniert.
      \end{itemize}
  \end{itemize}
\end{frame}
\begin{frame}
  \frametitle{Tree-Objekte}
  \begin{itemize}
    \item<1>Entsprechen Verzeichnissen im Dateisystem
    \item<1>Ebenfalls durch SHA1-Hashwert dargestellt
      \begin{itemize}
        \item unter Anderem abh�ngig von den enthaltenen Trees und Blobs
      \end{itemize}
    \item<1>Mindestens ein Blob- oder Tree-Objekt inklusive Namen zugeordnet
      \begin{itemize}
        \item Tree und Blob oder anderer Tree ergeben Datei-/Verzeichnisreferenz
      \end{itemize}
  \end{itemize}
\end{frame}
\begin{frame}
  \frametitle{Tree-Objekte Beispiel}
  \begin{center}
    \includegraphics{images/tree-example.png}
  \end{center}
\end{frame}
\begin{frame}
  \frametitle{Commit-Objekte}
  \begin{itemize}
    \item<1->Speichern Versionen der Daten
    \item<1->Beinhalten ein Tree-Objekt und einen oder mehrere Eltern-Commits
    \item<1->Dargestellt durch SHA1-Hashwert ($=$ Commit-ID)
      \begin{itemize}
        \item unter Anderem abh�ngig von den enthaltenen Trees, dem Autor sowie dem Datum und der Uhrzeit des Commits
      \end{itemize}
    \item<2->Commit-ID abh�ngig von den Trees $\Rightarrow$ abh�ngig von den Blobs $\Rightarrow$ abh�ngig von den Datei-Inhalten
    \item<3->Commit in Gesamtheit als \emph{Arbeitsschritt $X$} zum \emph{Zeitpunkt $Y$} von \emph{Person $Z$} weltweit eindeutig identifizierbar
    \item<4->(Sicherheits-)kopie aller Commits in jeder Arbeitskopie vorhanden
    \item<5->$\Rightarrow$ Sicherheitsvorteil (bsp. Linux Kernel)
  \end{itemize}
\end{frame}
\begin{frame}
  \frametitle{Commit-Objekte Beispiel}
  \begin{center}
    \includegraphics{images/commit-example.png}
  \end{center}
\end{frame}
\subsection[Remotes, Branches \& Tags]{Verteilung und Strukturierung der Historie}
\begin{frame}
  \frametitle{Remotes, Branches \& Tags}
\end{frame}
%
\section{GIT vs. SVN}
\begin{frame}
  \frametitle{GIT vs. SVN}
  \tableofcontents[currentsection]
\end{frame}
\subsection{R�ckblick zu SVN}
\begin{frame}
  \frametitle{R�ckblick zu SVN}
\end{frame}
\subsection{Workflow mit GIT}
\begin{frame}
  \frametitle{Workflow mit GIT}
\end{frame}
\subsection{Tabellarischer Vergleich}
\begin{frame}
  \frametitle{Tabellarischer Vergleich}
\end{frame}
%
\section{Beispiel mit zwei Entwicklern}
\begin{frame}
  \frametitle{Beispiel mit zwei Entwicklern}
  \tableofcontents[currentsection]
\end{frame}
\subsection{Bowling Game Kata}
\begin{frame}
  \frametitle{Bowling Game Kata}
\end{frame}
\end{document}
